% Part: intuitionistic-logic
% Chapter: tableaux
% Section: soundness

\documentclass[../../../include/open-logic-section]{subfiles}

\begin{document}

\olfileid{int}{tab}{sou}

\olsection{అంతఃప్రజ్ఞావాద \usetoken{టాబ్లోల}{tableau} నిర్దుష్టత}

\begin{explain}
  అంతఃప్రజ్ఞావాద \tetoken{టాబ్లోలు}{!!{tableau}s}
  నిర్దుష్టమైనవని చూపాలంటే, కిందివాటికి
  \[
  \sFmla{\True}{!B_1}[1], \dots, \sFmla{\True}{!B_n}[1], \sFmla{\False}{!A}[1]
  \]
  సంవృత \tetoken{టాబ్లో}{!!{tableau}} ఉంటే
  $!B_1, \dots, !B_n \Entails !A$ అని
  చూపాలి. ప్రతివిపర్యయం నిరూపించడం సులభం:
  ఏదో $\mModel{M}$, లోకం~$w$ వద్ద
  అన్ని $i=1$, \dots,~$n$లకు
  $\mSat{M}{!B_i}[w]$, కానీ
  $\mSat/{M}{!A}[w]$ అయితే ఏ
  \tetoken{టాబ్లోనూ}{!!{tableau}} సంవృతం
  చేయలేం.
  \sourcecorrection{OLTEINTTABSOU-001}{మూల ప్రతివిపర్యయ వాక్యం పూర్వాపేక్షలతో పాటు ముగింపూ సత్యమని వ్రాసింది; వ్యతిరేక నమూనా కావాలంటే ముగింపు అసత్యం కావాలి. ప్రతికూల సత్య సంకేతానికి మార్చాం.}
  అలాంటి ప్రతినమూనా \tetoken{టాబ్లో}{!!{tableau}}
  ప్రారంభ పరికల్పనలు సంతృప్తిపరచదగినవని
  చూపుతుంది. నిరూపణ వ్యూహం: ఒక
  \tetoken{టాబ్లో}{!!a{tableau}} శాఖపై ఉన్న
  అన్ని పూర్వసూచిక గల
  \tetoken{సూత్రాలు}{!!{formula}s}
  సంతృప్తిపరచదగినప్పుడు, ఏ నియమాన్ని
  వర్తింపజేసినా పొడిగించిన శాఖల్లో
  కనీసం ఒకటి సంతృప్తిపరచదగినదే
  అని చూపడం. సంవృత శాఖలు
  అసంతృప్తిపరచదగినవి కాబట్టి,
  పూర్వసూచిక గల \tetoken{సూత్రాల}{!!{formula}s}
  సంతృప్తిపరచదగిన సమితికి చెందిన
  ఏ \tetoken{టాబ్లోలోనైనా}{!!{tableau}}
  కనీసం ఒక వివృత శాఖ ఉండాలి.

  మోడల్ తర్క కేసులో ఈ వ్యూహాన్ని
  వర్తింపజేయాలంటే ``సంతృప్తిపరచదగినది''
  అనే నిర్వచనాన్ని సంబంధాలు,
  పూర్వసూచికలకు విస్తరించాలి.
  అది చేసిన తరువాత నిరూపణ నేరుగా సాగుతుంది.
\end{explain}

\begin{defn}
  $P$ పూర్వసూచికల సమితి అని, అంటే
  $P \subseteq (\PosInt)^*
  \setminus \{\emptyseq\}$ అని,
  $\mModel{M}$ ఒక నమూనా అని అనుకుందాం.
  $\sigma$, $\sigma.n$ రెండూ~$P$లో
  ఉన్నప్పుడల్లా $Rf(\sigma)f(\sigma.n)$
  అయితే, ప్రమేయం~$f\colon P \to W$ను
  $\mModel{M}$లో~$P$ యొక్క
  \emph{అర్థనిర్దేశం} అంటాం.

  పూర్వసూచికల అర్థనిర్దేశం~$P$కు
  సంబంధించి ఇలా నిర్వచించవచ్చు:
  \begin{enumerate}
  \item $\mSat{M}{!A}[f(\sigma)]$ అయితే,
    అప్పుడే $\mModel{M}$,
    $\sFmla{\True}{!A}[\sigma]$ను సంతృప్తిపరుస్తుంది.
  \item $\mSat/{M}{!A}[f(\sigma)]$ అయితే,
    అప్పుడే $\mModel{M}$,
    $\sFmla{\False}{!A}[\sigma]$ను సంతృప్తిపరుస్తుంది.
  \end{enumerate}
\end{defn}

$R$ స్వావర్తన, సంక్రామక; $\sigma.{*}$
అనేది $\sigma$, $\sigma.n_1$,
$\sigma.n_1.n_2$, \dots ను సూచిస్తుంది.
కాబట్టి $Rf(\sigma)f(\sigma.{*})$
కూడా ఉంటుందని గమనించండి.

\begin{defn}
  $\Gamma$ పూర్వసూచిక గల
  \tetoken{సూత్రాల}{!!{formula}s} సమితి,
  $P(\Gamma)$ అందులో కనిపించే
  పూర్వసూచికల సమితి అనుకుందాం.
  $f$ అనేది $\mModel{M}$లో~$P(\Gamma)$కు
  అర్థనిర్దేశమైతే, $\Gamma$లోని ప్రతి
  పూర్వసూచిక గల \tetoken{సూత్రాన్ని}{!!{formula}}
  $\mModel{M}$,~$f$కు సంబంధించి
  సంతృప్తిపరిచినప్పుడే, $\mModel{M}$
  అనేది~$f$కు సంబంధించి $\Gamma$ను
  సంతృప్తిపరుస్తుంది,
  $\mSat{M}{\Gamma}[f]$ అంటాం.
  $\mSat{M}{\Gamma}[f]$ అయ్యే నమూనా~$\mModel{M}$,
  $P(\Gamma)$కు అర్థనిర్దేశం~$f$ ఉంటే,
  అప్పుడే $\Gamma$ \emph{సంతృప్తిపరచదగినది}.
\end{defn}

\begin{prop}
  ఏదో \tetoken{సూత్రం}{!!{formula}}~$!A$,
  పూర్వసూచిక~$\sigma$కు
  $\Gamma$లో $\sFmla{\True}{!A}[\sigma]$,
  $\sFmla{\False}{!A}[\sigma.{*}]$
  రెండూ ఉన్నా, లేదా
  $\sFmla{\True}{\lfalse}[\sigma]$
  ఉన్నా, $\Gamma$ అసంతృప్తిపరచదగినది.
\end{prop}

\begin{proof}
  ఎల్లప్పుడూ $\mSat/{M}{\lfalse}[f(\sigma)]$
  కాబట్టి $\sFmla{\True}{\lfalse}$
  గల $\Gamma$ అసంతృప్తిపరచదగినది.

  $\mModel{M}$ నమూనా, $P(\Gamma)$కు
  అర్థనిర్దేశం~$f$ వద్ద
  $\mSat{M}{!A}[f(\sigma)]$ అయితే,
  \olref[sem][rel]{prop:true-monotonic}
  ప్రకారం, $Rf(\sigma)f(\sigma.{*})$
  కావడంతో
  $\mSat{M}{!A}[f(\sigma.{*})]$.
  \sourcecorrection{OLTEINTTABSOU-002}{మూల నిరూపణ వాక్యం వ్యాకరణపరంగా తెగిపోయి, ప్రాప్యత సంబంధంలో రెండో లోకానికి $f$ను వదిలింది; ఏకదిశ సత్యం వల్ల విస్తరించిన లోకంలోనే $!A$ సత్యమనే అవసరమైన ముగింపుకు సరిచేశాం.}
  కాబట్టి $\mSat{M}{!A}[f(\sigma)]$,
  $\mSat/{M}{!A}[f(\sigma.{*})]$
  రెండూ కలిసి ఉండలేవు.
\end{proof}

\begin{thm}[నిర్దుష్టత]
  \ollabel{thm:tableau-soundness}
  $\Gamma$కు సంవృత \tetoken{టాబ్లో}{!!{tableau}}
  ఉంటే, $\Gamma$ అసంతృప్తిపరచదగినది.
\end{thm}

\begin{proof}
\tetoken{ఒక టాబ్లో}{!!a{tableau}} శాఖపై ఉన్న
\tetoken{చిహ్నిత సూత్రాల}{!!{signed formula}s}
సమితి సంతృప్తిపరచదగినప్పుడే ఆ శాఖ
సంతృప్తిపరచదగినదని అందాం.
కనీసం ఒక సంతృప్తిపరచదగిన శాఖ ఉన్న
\tetoken{ఒక టాబ్లోను}{!!a{tableau}}
సంతృప్తిపరచదగినదని అందాం.

మనం చూపేది ఇది: సంతృప్తిపరచదగిన
\tetoken{టాబ్లోను}{!!{tableau}}
ఏదైనా నిగమన నియమంతో విస్తరిస్తే,
ఫలిత \tetoken{టాబ్లో}{!!{tableau}}
కూడా సంతృప్తిపరచదగినదే. దీనివల్ల
సిద్ధాంతం వస్తుంది: సంవృత
\tetoken{టాబ్లో}{!!{tableau}}
అనేది~$\Gamma$ పరికల్పనలు మాత్రమే గల
\tetoken{టాబ్లోకు}{!!{tableau}}
నిగమన నియమాలు వర్తింపజేయడం వల్ల
వస్తుంది. కాబట్టి $\Gamma$
సంతృప్తిపరచదగినదైతే, దానికి చెందిన
ప్రతి \tetoken{టాబ్లో}{!!{tableau}}
సంతృప్తిపరచదగినదే. కానీ సంవృత
\tetoken{టాబ్లో}{!!{tableau}}
సంతృప్తిపరచదగినది కాదు: దాని అన్ని
శాఖలు సంవృతమైనవి, సంవృత శాఖలు
అసంతృప్తిపరచదగినవి.

మనకు సంతృప్తిపరచదగిన
\tetoken{టాబ్లో}{!!{tableau}}, అంటే
కనీసం ఒక సంతృప్తిపరచదగిన శాఖ ఉన్న
\tetoken{ఒక టాబ్లో}{!!a{tableau}}
ఉందనుకుందాం. నిగమన నియమం ఒక
శాఖకు \tetoken{చిహ్నిత
సూత్రాలను}{!!{signed formula}s} చేరుస్తుంది,
లేదా శాఖను రెండుగా విడదీస్తుంది.
నియమం విస్తరించని సంతృప్తిపరచదగిన
శాఖ \tetoken{టాబ్లోలో}{!!{tableau}}
ఉంటే, విస్తరించిన
\tetoken{టాబ్లోలో}{!!{tableau}}
ఆ శాఖ సంతృప్తిపరచదగినదిగానే
ఉంటుంది; అందువల్ల విస్తరించిన
టాబ్లో కూడా సంతృప్తిపరచదగినదే.
కాబట్టి నియమాన్ని
సంతృప్తిపరచదగిన శాఖకు
వర్తింపజేసిన సందర్భాన్నే పరిశీలించాలి.

ఆ శాఖపై ఉన్న \tetoken{చిహ్నిత
సూత్రాల}{!!{signed formula}s} సమితి
$\Gamma$ అనుకుందాం; నియమం వర్తించే
\tetoken{చిహ్నిత సూత్రం}{!!{signed formula}}
$\sFmla{S}{!A}[\sigma] \in \Gamma$
అనుకుందాం. నియమం శాఖను
విడదీయకపోతే, $\Gamma$కు నియమపు
ఫలితాలను చేర్చిన విస్తరించిన శాఖ
ఇంకా సంతృప్తిపరచదగినదని చూపాలి.
శాఖను విడదీస్తే, ఫలిత రెండు
శాఖల్లో కనీసం ఒకటి
సంతృప్తిపరచదగినదని చూపాలి.
\begin{enumerate}
\item $\sFmla{\True}{\lnot !B}[\sigma] \in \Gamma$కు
  $\TRule{\True}{\lnot}$ వర్తింపజేసి
  శాఖను విస్తరిద్దాం. విస్తరించిన
  శాఖలోని \tetoken{చిహ్నిత
  సూత్రాలు}{!!{signed formula}s} $\Gamma \cup
  \{\sFmla{\False}{!B}[\sigma.{*}]\}$.
  $\mSat{M}{\Gamma}[f]$ అనుకుందాం.
  ముఖ్యంగా $\mSat{M}{\lnot !B}[f(\sigma)]$.
  అందువల్ల $Rf(\sigma)w$ అయ్యే ఏ $w$
  వద్దా $\mSat/{M}{!B}[w]$; అందులో
  $f(\sigma.{*})$ కూడా ఉంది.
  కాబట్టి $\mModel{M}$,~$f$కు
  సంబంధించి
  $\sFmla{\False}{!B}[\sigma.{*}]$ను
  సంతృప్తిపరుస్తుంది.
\item $\sFmla{\False}{\lnot !B}[\sigma] \in \Gamma$కు
  $\TRule{\False}{\lnot}$
  వర్తింపజేసి శాఖను విస్తరించండి:
  అభ్యాసం.
\item $\sFmla{\True}{!B \land !C}[\sigma] \in
  \Gamma$కు $\TRule{\True}{\land}$
  వర్తింపజేస్తే, శాఖపై రెండు కొత్త
  \tetoken{చిహ్నిత సూత్రాలు}{!!{signed formula}s}
  $\sFmla{\True}{!B}[\sigma]$,
  $\sFmla{\True}{!C}[\sigma]$ వస్తాయి.
  $\mSat{M}{\Gamma}[f]$ అనుకుందాం;
  ముఖ్యంగా $\mSat{M}{!B \land
    !C}[f(\sigma)]$. అప్పుడు
  $\mSat{M}{!B}[f(\sigma)]$,
  $\mSat{M}{!C}[f(\sigma)]$. అంటే
  $\mModel{M}$,~$f$కు సంబంధించి
  $\sFmla{\True}{!B}[\sigma]$,
  $\sFmla{\True}{!C}[\sigma]$
  రెండింటినీ సంతృప్తిపరుస్తుంది.
\item $\sFmla{\False}{!B \lor !C} \in \Gamma$కు
  $\TRule{\False}{\lor}$ వర్తింపజేసి
  శాఖను విస్తరించండి: అభ్యాసం.
\item $\sFmla{\False}{!B \lif !C}[\sigma] \in
  \Gamma$కు $\TRule{\False}{\lif}$
  వర్తింపజేస్తే, శాఖపై రెండు కొత్త
  \tetoken{చిహ్నిత సూత్రాలు}{!!{signed formula}s}
  $\sFmla{\True}{!B}[\sigma.n]$,
  $\sFmla{\False}{!C}[\sigma.n]$ వస్తాయి;
  ఇక్కడ $\sigma.n$ ఆ శాఖపై కొత్త
  పూర్వసూచిక, అంటే
  $\sigma.n \notin P(\Gamma)$.
  
  $\Gamma$ సంతృప్తిపరచదగినది కాబట్టి,
  $\mSat{M}{\Gamma}[f]$ అయ్యే
  $\mModel{M}$, $P(\Gamma)$కు
  అర్థనిర్దేశం~$f$ ఉన్నాయి; ముఖ్యంగా
  $\mSat/{M}{!B \lif !C}[f(\sigma)]$.
  $\Gamma \cup
  \{\sFmla{\True}{!B}[\sigma.n],
    \sFmla{\False}{!C}[\sigma.n]\}$
  సంతృప్తిపరచదగినదని చూపాలి.
  \sourcecorrection{OLTEINTTABSOU-003}{మూల నిరూపణలో కొత్త శాఖ లక్ష్యంగా అసత్య సోపాధికాన్ని కొత్త పూర్వసూచిక వద్ద మళ్లీ చేర్చింది; నియమం నిజంగా చేరుస్తున్న సత్య పూర్వపక్షం, అసత్య ఉత్తరపక్షం రెండింటిని లక్ష్యంగా సరిచేశాం.}
  ఇందుకోసం $P(\Gamma)
  \cup \{\sigma.n\}$కు అర్థనిర్దేశాన్ని
  ఇలా నిర్వచిద్దాం:
  
  $\mSat/{M}{!B \lif !C}[f(\sigma)]$
  కాబట్టి $Rf(\sigma)w$ అయ్యే, అలాగే
  $\mSat{M}{!B}[w]$,
  $\mSat/{M}{!C}[w]$ అయ్యే ఏదో
  $w \in W$ ఉంది. $f'$ను $f$లాగే
  తీసుకొని, $f'(\sigma.n)
  = w$గా మాత్రమే మార్చుదాం.
  $f'(\sigma) = f(\sigma)$,
  $Rf(\sigma)w$ కాబట్టి
  $Rf'(\sigma)f'(\sigma.n)$;
  అందువల్ల $f'$,
  $P(\Gamma) \cup \{\sigma.n\}$కు
  అర్థనిర్దేశం. స్పష్టంగా
  $\mSat{M}{!B}[f'(\sigma.n)]$,
  $\mSat/{M}{!C}[f'(\sigma.n)]$.
  అన్ని $\sigma' \in
  P(\Gamma)$లకు $f(\sigma') = f'(\sigma')$
  కాబట్టి $\mSat{M}{\Gamma}[f']$.
  అందువల్ల $\mModel{M}, f'$
  అనేవి $\Gamma \cup
  \{\sFmla{\True}{!B}[\sigma.n],
    \sFmla{\False}{!C}[\sigma.n]\}$
  ను సంతృప్తిపరుస్తాయి.
\end{enumerate}
ఇప్పుడు రెండు శాఖలు వచ్చే నిగమనాలను
పరిశీలిద్దాం.
\begin{enumerate}
\item $\sFmla{\False}{!B \land !C}[\sigma] \in \Gamma$కు
  $\TRule{\False}{\land}$ వర్తింపజేసి శాఖను
  విస్తరిస్తే, రెండు శాఖలు వస్తాయి:
  ఎడమ శాఖ $\sFmla{\False}{!B}[\sigma]$తో,
  కుడి శాఖ $\sFmla{\False}{!C}[\sigma]$తో
  కొనసాగుతాయి. $\mSat{M}{\Gamma}[f]$
  అనుకుందాం; ముఖ్యంగా
  $\mSat/{M}{!B \land !C}[f(\sigma)]$.
  అప్పుడు $\mSat/{M}{!B}[f(\sigma)]$
  లేదా $\mSat/{M}{!C}[f(\sigma)]$.
  మొదటి సందర్భంలో $\mModel{M}, f$,
  $\sFmla{\False}{!B}[\sigma]$ను
  సంతృప్తిపరుస్తాయి; అంటే ఎడమ శాఖ
  సంతృప్తిపరచదగినది. రెండో సందర్భంలో
  $\mModel{M}, f$,
  $\sFmla{\False}{!C}[\sigma]$ను
  సంతృప్తిపరుస్తాయి; అంటే కుడి శాఖ
  సంతృప్తిపరచదగినది.
\item $\sFmla{\True}{!B \lor !C}[\sigma] \in \Gamma$కు
    $\TRule{\True}{\lor}$ వర్తింపజేసి
    శాఖను విస్తరించండి: అభ్యాసం.
\item $\sFmla{\True}{!B \lif !C}[\sigma] \in \Gamma$కు
    $\TRule{\True}{\lif}$ వర్తింపజేసి
    శాఖను విస్తరించండి: అభ్యాసం.
\end{enumerate}
\end{proof}

\begin{prob}
\olref[int][tab][sou]{thm:tableau-soundness}
నిరూపణను పూర్తి చేయండి.
\end{prob}

\begin{cor}
\ollabel{cor:entailment-soundness}
$\Gamma \Proves !A$ అయితే
$\Gamma \Entails !A$.
\end{cor}

\begin{proof}
  $\Gamma \Proves !A$ అయితే, ఏవో
  $!B_1$, \dots, $!B_n \in
  \Gamma$లకు, $\Delta =
  \{\sFmla{\False}{!A}[1], \sFmla{\True}{!B_1}[1],
  \dots, \sFmla{\True}{!B_n}[1]\}$కు
  సంవృత \tetoken{టాబ్లో}{!!{tableau}}
  ఉంది. $\Gamma \Entails !A$ అని
  చూపాలి. అలా కాదనుకుందాం. అప్పుడు
  ఏదో $\mModel{M}$, $w$కు,
  $i=1$, \dots,~$n$ అన్నింటికీ
  $\mSat{M}{!B_i}[w]$, కానీ
  $\mSat/{M}{!A}[w]$.
  $f(1) = w$గా తీసుకుందాం. అప్పుడు
  $f$, $P(\Delta)$కు $\mModel{M}$లో
  అర్థనిర్దేశం; $\mModel{M}$,~$f$కు
  సంబంధించి $\Delta$ను
  సంతృప్తిపరుస్తుంది. కానీ
  \olref{thm:tableau-soundness}
  ప్రకారం $\Delta$కు సంవృత
  \tetoken{టాబ్లో}{!!{tableau}}
  ఉన్నందున అది అసంతృప్తిపరచదగినది;
  ఇది విరోధం. కాబట్టి
  $\Gamma \Entails !A$ కావాలి.
  \sourcecorrection{OLTEINTTABSOU-004}{మూల చివరి వాక్యం ఇప్పటికే పూర్వాపేక్షగా తీసుకున్న $\Gamma \Proves !A$నే మళ్లీ ముగింపుగా పేర్కొంది; నిరూపించాల్సిన అర్థపరమైన పర్యవసానం $\Gamma \Entails !A$గా సరిచేశాం.}
\end{proof}

\begin{cor}
\ollabel{cor:weak-soundness}
$\Proves !A$ అయితే
$!A$ అన్ని నమూనాల్లో సత్యం.
\end{cor}

\end{document}
